\section{The Morse-number question}\label{sec:introduction}
Arnold's fixed-point conjectures compare Hamiltonian fixed points with
critical points of functions on the underlying symplectic manifold
\cite[Section~2]{Arnold1986}. For a sufficiently $C^2$-small autonomous Hamiltonian,
the fixed points of its time-one map are precisely its critical points.
The question is how much of this finite-dimensional lower bound survives
for an arbitrary Hamiltonian path. We address the formulation using the
ordinary Morse number.

Let $(M,\omega)$ be a closed connected symplectic manifold and let
$H\in C^\infty((\R/\Z)\times M,\R)$. We use the convention
\[
 \omega(X_H(t,\cdot),\cdot)=dH_t,
 \qquad \phi_H^0=\mathrm{id}_M
\]
for the Hamiltonian vector field and its flow. A fixed point $x$ of
$\phi_H^1$ is \emph{nondegenerate} if
$1\notin\Spec(d_x\phi_H^1)$. Write $\Fix_0(\phi_H^1;H)$ for the fixed
points whose loops $t\mapsto\phi_H^t(x)$ are contractible in $M$.
Contractibility here refers to this Hamiltonian path.

The \emph{ordinary Morse number} is
\[
 \Morse(M)=\min\{\#\Crit(f):f\in C^\infty(M,\R)\text{ is Morse}\}.
\]
The Morse-number form of the Arnold conjecture asks whether
$\#\Fix_0(\phi_H^1;H)\ge\Morse(M)$ whenever the time-one map is
nondegenerate \cite[Conjecture~1.1]{Golovko2020}.
We construct examples in which even the total number of fixed points is
smaller than $\Morse(M)$, and every fixed-point trajectory is contractible.

\begin{theorem}\label{thm:main}
There exist a closed connected symplectic manifold $(M,\omega)$ of real
dimension twelve and a smooth one-periodic Hamiltonian $H$ such that every
fixed point of $\phi_H^1$ is nondegenerate and
\[
\#\Fix_0(\phi_H^1;H)=\#\Fix(\phi_H^1)
                   =\chi(M)+128<\Morse(M).
\]
\end{theorem}

The same construction gives a quantitative separation in fixed dimension.
\begin{corollary}\label{cor:deficit}
For every real number $D>0$, there is an example as in
Theorem~\ref{thm:main} satisfying
\[
 \Morse(M)-\#\Fix(\phi_H^1)>D.
\]
The manifold is allowed to depend on $D$.
\end{corollary}

\subsection{From homological bounds to ordinary Morse theory}
Conley and Zehnder proved the nondegenerate bound on the standard
symplectic torus, where the ordinary Morse number equals the total Betti
number~\cite{ConleyZehnder1983}.
Floer's pseudoholomorphic-curve approach established homological Morse
inequalities for Hamiltonian fixed points, including the monotone
setting~\cite{Floer1989}. Virtual-moduli constructions of Fukaya--Ono,
Liu--Tian and Ruan extended the rational Betti-number bound to general
closed symplectic manifolds
\cite{FukayaOno1999,LiuTian1998,Ruan1999}.
The integral bounds of Bai--Xu retain torsion information as well
\cite[Theorem~A]{BaiXu2022}, \cite[Corollary~1.2]{BaiXu2026}.
These homological bounds do not identify the number of generators of a
Floer complex with the minimum number of critical points of an ordinary
Morse function.

Stabilization explains one source of this distinction. The
\emph{stable Morse number} minimizes the number of critical points of
Morse functions $F:M\times\R^r\to\R$ that are almost quadratic at
infinity, allowing the integer $r\geq0$ to vary. Here this means that $\|dF-dQ\|$ is
uniformly bounded in a product metric for some nondegenerate quadratic
form $Q$ on $\R^r$~\cite[Section~2.1]{DRG2017}.
Arnold already considered a related stable critical-point minimum and
asked whether it equals the minimum on the original manifold
\cite[Section~5]{Arnold1986}.
Eliashberg and Gromov developed the relationship with simple chain
homotopy over the fundamental-group ring and augmentation-ideal
generation~\cite[Sections~0.2.2 and~0.2.5]{EliashbergGromov1998}.
Damian proved that ordinary and stable Morse numbers of closed manifolds
can differ~\cite{Damian2002}. A relative realization theorem after
sufficiently large stabilization is given in
\cite[Proposition~2.9]{DRG2018}.

The algebraic distinction used here is classical and particularly sharp
for direct powers of nontrivial perfect finite groups. Cossey, Gruenberg and
Kov\'acs proved that the augmentation-ideal generator number remains
constant under these direct powers, while the group generator number
is unbounded~\cite[Theorems~1--2]{CGK1974}; see also
\cite[Sections~1--2]{GruenbergLinnell}.
Fundamental-group refinements of Floer theory reflect this distinction:
Ono--Pajitnov obtain augmentation-ideal bounds under weak monotonicity
\cite[Theorem~5.7]{OnoPajitnov2016}, whereas, in the symplectically
aspherical or monotone setting, Barraud's reconstruction of group
generators counts Floer trajectories with multiplicities
\cite[Theorem~1.3 and Remark~3]{Barraud2018}.
The trajectory multiplicities are part of the bound; these results do
not yield the unrestricted ordinary Morse-number inequality.

For closed symplectic manifolds on which both $[\omega]$ and $c_1(TM)$
vanish on $\pi_2(M)$, Dimitroglou Rizell and Golovko prove the stable
Morse number as a lower bound for nondegenerate Hamiltonian fixed
points~\cite[Corollary~1.3]{DRG2017}.
The sphere factors in our examples have positive symplectic area, so
this hypothesis does not hold. The construction below instead realizes
a short relative Morse complex in exactly eight additional variables,
with the function equal to a prescribed quadratic form outside a compact
set. Its small-time return map is then inserted into a closed product
of spheres. The relative handle realization and the exclusion of all
additional fixed points are the two geometric steps needed to turn the
ordinary/stable distinction into the stated closed Hamiltonian example.

The integral homology of our examples is torsion-free and concentrated
in even degrees, with total rank $\chi(M)$, as verified in
Section~\ref{sec:counting}. Their fixed-point count therefore exceeds
the homological quantities above.

Ma's preprint asserts that the minimum number of fixed points of a
nondegenerate Hamiltonian diffeomorphism always equals the ordinary
Morse number~\cite[Theorem~1.5]{Ma2013}. Theorem~\ref{thm:main}
contradicts that unrestricted assertion.

\subsection{The construction and its two counts}
For a finitely generated group $G$, let $d(G)$ be the least number of group
generators. Set $b=2$ and consider the finite groups $G=A_5^k$, where
$A_5$ has order $60$. Two elementary algebraic facts will drive the proof:
$d(A_5^k)$ is unbounded as $k$ increases, whereas the right augmentation
ideal
\[
 I_G=\ker(\aug:\Z[G]\longrightarrow\Z)
\]
has a generating list of length $b$, independent of $k$.
The distinction is between generation of a group and generation of a
module over its group ring.

Gompf's existence theorem supplies a closed connected symplectic
four-manifold $X$ with $\pi_1(X)=G$~\cite[Theorem~3.1]{Gompf1992};
see also \cite{Gompf1995}.
We construct a Morse function on $X\times\R^8$ with precisely
$\chi(X)+4b$ critical points, equal outside a compact set to the quadratic
form of index four
\[
 Q(u,v)=\frac{|v|^2-|u|^2}{2},\qquad u,v\in\R^4.
\]
The eight variables are fixed throughout: the proof performs handle moves
on one twelve-dimensional cobordism, rather than appealing to a
stabilization of unspecified dimension.

The closed symplectic manifold is
$M=X\times(S^2)^4$. Sixteen disjoint product-cap neighborhoods correspond
to the north/south choices on the four sphere factors. We place the same
small-time return map of the constructed Morse function in each cap.
The main gluing issue is to match the local \emph{endpoint} with product
rotations at both poles while excluding fixed points outside the caps.
The resulting smooth one-periodic Hamiltonian satisfies
\begin{equation}\label{eq:overview-fixed-count}
 \#\Fix(\phi_H^1)=16\bigl(\chi(X)+4b\bigr)
                   =\chi(M)+128.
\end{equation}

On the other hand, the fundamental group forces every Morse function on
$M$ to have at least $d(G)$ critical points in each of indices one and
eleven. Euler's identity gives
\begin{equation}\label{eq:overview-morse-bound}
 \Morse(M)\ge\chi(M)+4d(G).
\end{equation}
Since $d(A_5^k)$ is unbounded while $b$ is fixed, these two estimates
give the strict inequality and the arbitrary deficits. Figure~\ref{fig:construction}
separates the two roles of the fundamental group.

\begin{figure}[ht]
\centering
\begin{tikzpicture}[
 every node/.style={align=center,font=\small},
 box/.style={draw,inner sep=6pt,text width=58mm},
 >={Stealth}]
\node[box,text width=90mm] (group) at (0,0)
{$G=A_5^k$\\$I_G$ has two generators; $d(G)$ is unbounded};
\node[box] (morse) at (-3.4,-1.8)
 {Relative handle cancellation\\$X^4\times\R^8$\\$\chi(X)+8$ critical points};
\node[box] (bound) at (3.4,-1.8)
 {Ordinary Morse lower bound\\$M^{12}=X\times(S^2)^4$\\$\Morse(M)\ge\chi(M)+4d(G)$};
\node[box] (closed) at (-3.4,-3.6)
 {Sixteen cap insertions\\$\#\Fix(\phi_H^1)=\chi(M)+128$};
\node[box] (gap) at (3.4,-3.6)
 {Unbounded deficit\\$\Morse(M)-\#\Fix(\phi_H^1)$};
\draw[->] (group)--(morse);
\draw[->] (group)--(bound);
\draw[->] (morse)--(closed);
\draw[->] (bound)--(gap);
\draw[->] (closed)--(gap);
\end{tikzpicture}
\caption{The augmentation-ideal bound controls the constructed fixed-point
count. The group generator bound supplies the larger ordinary Morse
lower bound on the same closed manifold.}
\label{fig:construction}
\end{figure}

\subsection{Proof organization}
Section~\ref{sec:algebra} proves the two group-ring facts and the exact
row operation needed for geometric cancellation.
Section~\ref{sec:relative-model} constructs the relative Morse model and
its handle complex. Section~\ref{sec:cancellation} reduces the handle
counts while preserving the quadratic collars.
Section~\ref{sec:hamiltonian} realizes the local map on the closed product.
Section~\ref{sec:counting} compares the two counts and proves the theorem
and its deficit corollary.
